\documentclass[10pt,a4paper]{amsart}
\usepackage[margin=19mm]{geometry}
\usepackage{amsmath,amssymb,mathtools}
\usepackage[scr=rsfs,cal=euler]{mathalfa}
\usepackage{xfrac}
\usepackage[unicode,hidelinks,bookmarks=false]{hyperref}
\newcommand{\ie}{i.e.\ }
\newcommand{\cf}{cf.\ }
\newcommand{\resp}{respectively\ }
\newcommand{\Subsection}{\subsection}
\providecommand{\nobreakdash}{\nobreak\hbox{-}}
\setlength{\emergencystretch}{2em}
\multlinegap0pt
\usepackage{mathtools}
\usepackage{amssymb}
\usepackage{bbm}
\usepackage[shortlabels]{enumitem}
\usepackage[normalem]{ulem}
\usepackage{caption}
\usepackage{tikz}
\usepackage{float}
\newtheorem{lemma}{Lemma}
\title{Quantitative longest-run laws for partial quotients}
\author{Ying Wai Lee}
\date{Source: arXiv:2602.11462v1 (CC BY 4.0)}
\begin{document}
\maketitle

\noindent\emph{Source and licence.} Quantitative longest-run laws for partial quotients. Version arXiv:2602.11462v1 (CC BY 4.0), distributed by arXiv under the Creative Commons Attribution 4.0 International licence, http://creativecommons.org/licenses/by/4.0/, as recorded in the arXiv licence metadata for this exact version and shown on the abstract page https://arxiv.org/abs/2602.11462v1. Full licence notice: \texttt{licenses/arxiv-2602.11462-CC-BY-4.0.txt}.\par\smallskip

\noindent\textit{Editorial scope.} Counted unit: the article's lemma 2 (source0.tex lines 433--438). The article's complete original proof of it is delivered with it (source0.tex lines 439--456, one \texttt{proof} environment of 18 lines). No mathematics is written, repaired or re-derived: the statement and the proof are the article's own lines, in the article's own notation, and no retained line is substituted or re-flowed.  Retained uncounted as the source-local context the proof needs: lines 363--365; lines 423--426.  Nothing was omitted from any retained span, so the package carries no omission record.  No retained line contains a citation, so the wrapper carries no bibliography and no bibliography entry.  No retained line contains a figure environment, an image-inclusion command or drawing code, so no image is delivered and the package declares no asset dependency.  Editorial formatting only, in addition: an AMS \texttt{amsart} preamble that restates the article's own declarations and macros used by the retained lines, namely the environments \texttt{lemma} on separate counters, together with the standard packages those lines need.  The retained environments print in this document as Lemma 1 in order of appearance, whereas the article prints the same items with the numbering of its own full document.  The complete native source of the article as distributed in the arXiv e-print for this exact version is bundled unchanged as \texttt{sources/194475/source0.tex}.
\begin{align}\label{eq: 16}
    \mu{(A)}\coloneqq\frac{1}{\log 2}\int_A\frac{\mathrm{d}x}{1+x},
\end{align}

\begin{align}\label{eq: 17}
    \operatorname{diam}{\Delta{(d_1,\dots,d_k)}}
    =\frac{1}{q_k(q_k+q_{k-1})}.
\end{align}

\begin{lemma}\label{lem: assumption 2}
For any $\lambda\in\mathbb{N}$, there exist $c_-,c_+>0$ such that for any $k\in\mathbb{N}$,
\begin{align*}
    c_-\,\tau^{-2k}\leq \mu{(\Delta_k{(\lambda)})}\leq c_+\,\tau^{-2k}.
\end{align*}
\end{lemma}

\begin{proof}
Pick any $k\in\mathbb{N}$. The denominator $q_k$ of the $k$-th convergent of $[\lambda,\lambda,\ldots]$ satisfies:
\begin{align*}
    \frac{\lambda}{\sqrt{\lambda^2+4}}\tau^{k}\leq q_{k}=\frac{\tau^{k+1}-(-\tau)^{-(k+1)}}{\sqrt{\lambda^2+4}}\leq \tau^{k}.
\end{align*}
By~\eqref{eq: 17},
\begin{align*}
    \frac{\tau}{\tau+1} \tau^{-2k}
    \leq \operatorname{diam}{\Delta_k{(\lambda)}}
    =\frac{1}{q_k(q_k+q_{k-1})}
    \leq\frac{(\lambda^2+4)\tau}{(\tau+1)\lambda^2}\tau^{-2k}
\end{align*}
By~\eqref{eq: 16},
\begin{align*}
     \frac{1}{2\log{2}}\frac{\tau}{\tau+1}\tau^{-2k}\leq \mu{(\Delta_k{(\lambda)})}\leq \frac{(\lambda^2+4)\tau}{(\tau+1)\lambda^2\log{2}}\tau^{-2k}.
\end{align*}
Lemma~\ref{lem: assumption 2} is proved.
\end{proof}

\end{document}
